\documentclass[11pt]{article}

\usepackage[margin=1in]{geometry}
\usepackage[T1]{fontenc}
\usepackage[utf8]{inputenc}

\title{Quad Mini Mesaj ve Bridge Notlari}
\author{}
\date{\today}

\begin{document}
\maketitle

\section*{Amac}
Bu dokuman, sistemde hangi mesajin nereden geldigini, hangi bridge node'un ne yaptigini ve \texttt{sim} ile \texttt{real} calisma akislarinin nasil kullanildigini kisa ve net sekilde ozetler.

\section*{Ana Mesajlar}
\begin{itemize}
  \item \texttt{legged\_msgs/JointControlData}: MPC/WBC tarafinin urettigi ust seviye eklem komutu.
  \item \texttt{stm2ros/JointControlState}: alt seviye donanim tarafina giden komut.
  \item \texttt{legged\_msgs/SimulatorSensorData}: \texttt{stateEstimate true} iken kullanilan sensor tabanli durum mesaji.
  \item \texttt{legged\_msgs/SimulatorStateData}: eski tam-durum mesaji. Sadece \texttt{stateEstimate false} akisinda kullanilir.
  \item \texttt{/odom}: tahmini govde pozu.
  \item \texttt{/joint\_states}: RViz ve \texttt{robot\_state\_publisher} icin standart eklem durumu.
  \item \texttt{/odom\_path}: \texttt{/odom}'dan uretilen iz mesaji.
\end{itemize}

\section*{Mesaj Icerikleri}
\subsection*{\texttt{JointControlData}}
\begin{itemize}
  \item \texttt{joint\_torque}
  \item \texttt{joint\_position}
  \item \texttt{joint\_velocity}
  \item \texttt{actuator\_mode}
\end{itemize}

Bu mesaj kontrol tarafindan uretilir.

\subsection*{\texttt{JointControlState}}
\begin{itemize}
  \item \texttt{joint\_names[12]}
  \item \texttt{joint\_position[12]}
  \item \texttt{joint\_torque[12]}
  \item \texttt{kp}
  \item \texttt{kd}
\end{itemize}

Bu mesaj alt seviye motor tarafina giden komuttur.

\subsection*{\texttt{SimulatorSensorData}}
\begin{itemize}
  \item \texttt{simulation\_time}
  \item \texttt{imu\_quat\_values}
  \item \texttt{imu\_angvel\_values}
  \item \texttt{imu\_linacc\_values}
  \item \texttt{joint\_position\_values}
  \item \texttt{joint\_velocity\_values}
  \item \texttt{contact\_flags}
\end{itemize}

Bu mesaj estimator, RViz ve tahmini durum akisinin ana mesajidir.

\section*{Temel Akis Grafigi}
\subsection*{1. Kontrol Komutu Akisi}
\begin{center}
\fbox{\texttt{MPC / WBC}}
$\rightarrow$
\fbox{\texttt{joint\_control\_data}}
$\rightarrow$
\fbox{\texttt{sim2real.py}}
$\rightarrow$
\fbox{\texttt{stm2ros/JointControlState}}
\end{center}

\subsection*{2. RViz Icin Tahmini Durum Akisi}
\begin{center}
\fbox{\texttt{simulator\_sensor\_data}}
$\rightarrow$
\fbox{\texttt{sensor\_to\_joint\_state.py}}
$\rightarrow$
\fbox{\texttt{/joint\_states}}
$\rightarrow$
\fbox{\texttt{robot\_state\_publisher}}
\end{center}

\begin{center}
\fbox{\texttt{/odom}}
$\rightarrow$
\fbox{\texttt{odom\_to\_tf.py}}
$\rightarrow$
\fbox{\texttt{TF: odom $\rightarrow$ base}}
\end{center}

\begin{center}
\fbox{\texttt{/odom}}
$\rightarrow$
\fbox{\texttt{odom\_to\_path.py}}
$\rightarrow$
\fbox{\texttt{/odom\_path}}
$\rightarrow$
\fbox{\texttt{RViz izi}}
\end{center}

\subsection*{3. Sim ve Real Icin Ortak Kural}
\begin{center}
\begin{tabular}{c}
\fbox{\texttt{./run.sh quad\_mini\_real sim estimated ...}} \\
$\downarrow$ \\
\fbox{\texttt{simulator\_sensor\_data}} \\
\\
\fbox{\texttt{./run.sh quad\_mini\_real real estimated ...}} \\
$\downarrow$ \\
\fbox{\texttt{simulator\_sensor\_data}}
\end{tabular}
\end{center}

Yani \texttt{stateEstimate true} iken RViz icin ortak kural sudur:
\begin{itemize}
  \item eklem bilgisi \texttt{simulator\_sensor\_data}'dan gelir
  \item govde hareketi \texttt{/odom}'dan gelir
  \item iz goruntusu \texttt{/odom\_path}'dan gelir
\end{itemize}

\section*{Bridge Scriptleri}
\subsection*{\texttt{sim2real.py}}
\begin{itemize}
  \item \texttt{JointControlData} alir
  \item \texttt{stm2ros/JointControlState} yayimlar
  \item canli \texttt{kp/kd} oran ayari vardir
  \item motor kapama ve feedforward torque kapama secenekleri vardir
\end{itemize}

\subsection*{\texttt{sensor\_to\_joint\_state.py}}
\begin{itemize}
  \item \texttt{simulator\_sensor\_data} icindeki eklem acilarini okur
  \item bunu \texttt{/joint\_states} formatina cevirir
  \item RViz'de bacaklarin gorunmesini saglar
\end{itemize}

\subsection*{\texttt{odom\_to\_tf.py}}
\begin{itemize}
  \item \texttt{/odom} mesajindan \texttt{odom $\rightarrow$ base} TF'i uretir
  \item sifira yakin hatali zaman damgasi gelirse mevcut zamani kullanir
  \item boylece RViz'de govde ve bacak TF zinciri duzgun oturur
\end{itemize}

\subsection*{\texttt{odom\_to\_path.py}}
\begin{itemize}
  \item \texttt{/odom}'dan \texttt{/odom\_path} uretir
  \item iz zaman tabanli degil, nokta sayisi tabanlidir
  \item ana parametreler: \texttt{max\_length} ve \texttt{min\_position\_delta}
\end{itemize}

\section*{Hardware Bridge Detayi}
Bu kisim \texttt{real\_robot\_bridge} icindeki \texttt{HardwareBackend} tarafini anlatir. Bu notta kullanilan varsayim sudur: sizin sisteminizde \texttt{htdw\_joint\_cmd} motorlardan gelen geri besleme tarafinda kullaniliyor. Yani burada \texttt{htdw\_joint\_cmd} bir feedback konusu olarak dusunulur. Gercek robot kullaniminda Orin tarafindan ayrica \texttt{SimulatorSensorData} gondermek zorunlu degildir. En temiz akista Orin ham ROS mesajlarini yayimlar, bridge bu mesajlardan \texttt{SimulatorSensorData} uretir ve gerekiyorsa temas kestirimi de bridge icinde yapilir.

\subsection*{Orin Tarafindan Yayimlanmasi Gerekenler}
\begin{itemize}
  \item \texttt{htdw\_joint\_cmd}: motorlardan gelen eklem geri beslemesi
  \item \texttt{sensor\_msgs/Imu} konu basligi: varsayilan konu \texttt{imu/data}
  \item \texttt{nav\_msgs/Odometry} konu basligi: istege bagli, parametre ile verilir
\end{itemize}

Yani bu notta bridge'in eklem geri besleme kaynagi dogrudan \texttt{htdw\_joint\_cmd} olarak dusunulur. Bridge'in ihtiyaci olan sey, bu feedback mesaji icinde eklem pozisyonu, hizi ve tork bilgilerinin bulunmasidir.

\subsection*{\texttt{htdw\_joint\_cmd} Feedback Icindeki Gerekli Alanlar}
\begin{itemize}
  \item eklem isimleri
  \item 12 eklem pozisyonu
  \item 12 eklem hizi
  \item 12 eklem torku
\end{itemize}

Tork bilgisi burada ozellikle onemlidir. Cunku bridge icindeki temas kestirimi bu feedback uzerinden gelen eklem torklarini kullanir.

\subsection*{Imu Icindeki Gerekli Alanlar}
\begin{itemize}
  \item \texttt{orientation}
  \item \texttt{angular\_velocity}
  \item \texttt{linear\_acceleration}
\end{itemize}

Bu alanlar \texttt{SimulatorSensorData} icindeki IMU kismini doldurur.

\subsection*{Odometry Icindeki Gerekli Alanlar}
\begin{itemize}
  \item \texttt{pose.pose.position}
  \item \texttt{twist.twist.linear}
\end{itemize}

Eger \texttt{odom} varsa bridge taban pozisyonu ve lineer hizi buradan alir. Eger yoksa varsayilan taban yuksekligi ve sifir hiz kullanilir.

\subsection*{Bridge Iceride Neyi Doldurur}
\texttt{HardwareBackend} ham mesajlari alip iki farkli yapiyi doldurur:
\begin{itemize}
  \item \texttt{data.state}
  \item \texttt{data.sensor}
\end{itemize}

Buradaki mantik sunudur:
\begin{itemize}
  \item \texttt{data.sensor}: estimator ve RViz icin kullanilan sensor tarafi
  \item \texttt{data.state}: bridge icindeki bazi legacy ve yardimci hesaplar icin kullanilan durum tarafi
\end{itemize}

\subsection*{Bridge'in Urettigi Sensor Mesaji}
\texttt{legged\_msgs/SimulatorSensorData} su alanlardan olusur:
\begin{itemize}
  \item \texttt{simulation\_time}
  \item \texttt{imu\_quat\_values}
  \item \texttt{imu\_angvel\_values}
  \item \texttt{imu\_linacc\_values}
  \item \texttt{joint\_position\_values}
  \item \texttt{joint\_velocity\_values}
  \item \texttt{contact\_flags}
\end{itemize}

Yani bridge su isi yapar:
\begin{center}
\fbox{\texttt{htdw\_joint\_cmd + Imu + Odom}}
$\rightarrow$
\fbox{\texttt{HardwareBackend}}
$\rightarrow$
\fbox{\texttt{SimulatorSensorData}}
\end{center}

\subsection*{Temas Kestirimi Nerede Yapilir}
Temas kestirimi bridge icinde yapilabilir ve mevcut yapi buna uygundur.

Akis su sekildedir:
\begin{center}
\fbox{\texttt{htdw\_joint\_cmd feedback torku}}
$\rightarrow$
\fbox{\texttt{data.state.joint\_torque\_values}}
$\rightarrow$
\fbox{\texttt{ContactEstimator}}
$\rightarrow$
\fbox{\texttt{contact\_flags}}
$\rightarrow$
\fbox{\texttt{SimulatorSensorData.contact\_flags}}
\end{center}

Bu nedenle gercek robotta temas kestirimi istiyorsan:
\begin{itemize}
  \item \texttt{contactSource = estimated} secilmelidir
  \item \texttt{htdw\_joint\_cmd} geri beslemesi gercek tork bilgisini tasimalidir
\end{itemize}

\subsection*{Orin'den Ayrica SimulatorSensorData Gondermek Gerekir mi}
Hayir, mevcut yapida buna mecbur degilsin.

En temiz kullanim:
\begin{itemize}
  \item Orin ham verileri yayimlar: \texttt{htdw\_joint\_cmd} feedback, \texttt{Imu}, gerekirse \texttt{Odometry}
  \item bridge bunlardan \texttt{SimulatorSensorData} olusturur
  \item bridge temas bayraklarini ekler
  \item estimator, RViz ve diger taraflar son durumda \texttt{simulator\_sensor\_data} kullanir
\end{itemize}

Bu mimari, bridge'i tek sorumlu nokta yapar ve sim ile real arasinda davranisi daha tutarli hale getirir.

\subsection*{Geri Besleme ve Komut Konulari}
Bu raporda kullanilan isimlendirme ile:
\begin{itemize}
  \item \texttt{htdw\_joint\_cmd}: motorlardan gelen eklem geri beslemesi
  \item ayri bir cikis komut konusu: motorlara giden komut tarafidir
\end{itemize}

Yani ayni konuyu hem feedback hem command icin kullanmamak gerekir. Bridge tarafinda mantik su olmali:
\begin{center}
\fbox{\texttt{htdw\_joint\_cmd (feedback)}}
$\rightarrow$
\fbox{\texttt{HardwareBackend::read}}
\end{center}

\begin{center}
\fbox{\texttt{JointControlData}}
$\rightarrow$
\fbox{\texttt{HardwareBackend::writeCommand}}
$\rightarrow$
\fbox{\texttt{ayri bir command topic}}
\end{center}

Bu ayri command topic, sistemde ne kullaniliyorsa o olabilir. Ornegin \texttt{joint\_cmd}, \texttt{motor\_command} veya baska bir ad.

\subsection*{Bridge'e Gelen Ust Seviye Komut}
Bridge sadece sensor tarafi degil, komut tarafini da yonetir. Ust seviye kontrol komutu:
\begin{itemize}
  \item \texttt{joint\_position}
  \item \texttt{joint\_velocity}
  \item \texttt{joint\_torque}
  \item \texttt{actuator\_mode}
\end{itemize}

\texttt{JointControlData} icinde gelir ve bridge bunu alt seviye komut konusuna yazar:
\begin{center}
\fbox{\texttt{JointControlData}}
$\rightarrow$
\fbox{\texttt{HardwareBackend::writeCommand}}
$\rightarrow$
\fbox{\texttt{alt seviye command topic}}
\end{center}

\section*{Nasil Kullanilir}
\subsection*{1. Rebuild}
\texttt{run.sh} install altindaki launch ve RViz dosyalarini kullandigi icin degisiklikten sonra rebuild gerekir:
\begin{verbatim}
./rebuild_quick.sh
\end{verbatim}

\subsection*{2. Simulasyon}
\begin{verbatim}
./run.sh quad_mini_real sim estimated debug rviz gui
\end{verbatim}

Bu durumda:
\begin{itemize}
  \item \texttt{simulator\_sensor\_data} MuJoCo tarafindan gelir
  \item RViz eklemleri buradan cizer
  \item \texttt{/odom} tahmini govde hareketini tasir
\end{itemize}

\subsection*{3. Gercek Robot}
\begin{verbatim}
./run.sh quad_mini_real real estimated debug rviz gui
\end{verbatim}

Bu durumda:
\begin{itemize}
  \item gercek robot bridge'i \texttt{simulator\_sensor\_data} yayimlar
  \item RViz yine ayni estimator akisindan beslenir
  \item yani sim ve real icin RViz mantigi aynidir
\end{itemize}

\subsection*{4. Debug}
\texttt{debug} aciksa:
\begin{itemize}
  \item \texttt{stateEstimate false} iken legacy state loglari basilir
  \item \texttt{stateEstimate true} iken sensor mesaji loglari basilir
\end{itemize}

\section*{RViz Tarafi}
Kalman RViz ayarlari su sekildedir:
\begin{itemize}
  \item \texttt{Fixed Frame = odom}
  \item kamera hedefi \texttt{base}
  \item TF gostergesi varsayilan olarak kapali
  \item eski odometry eksen gostergesi kapali
  \item bunun yerine \texttt{/odom\_path} uzerinden bir iz cizilir
\end{itemize}

\section*{Kisa Notlar}
\begin{itemize}
  \item \texttt{stateEstimate true} ise ana state mesaji \texttt{simulator\_sensor\_data}'dir.
  \item \texttt{sim2real.py}, \texttt{run.sh} ile otomatik baslamaz; gerekiyorsa ayri calistirilir.
  \item Iz goruntusunun uzunlugu esas olarak \texttt{odom\_to\_path.py} icindeki \texttt{max\_length} ile belirlenir.
  \item Iz yogunlugu ise \texttt{min\_position\_delta} ile ayarlanir.
  \item Gercek robotta temas kestirimi icin \texttt{htdw\_joint\_cmd} geri beslemesi icinde tork bilgisi gereklidir.
  \item Mevcut yapiyla Orin'den dogrudan \texttt{SimulatorSensorData} gondermek yerine ham sensor mesajlarini gondermek daha temizdir.
  \item Bu notta \texttt{htdw\_joint\_cmd} geri besleme konusu olarak kabul edilmistir.
  \item Bu durumda motorlara giden komut icin ayri bir topic kullanilmalidir.
\end{itemize}

\end{document}
